\documentclass[11pt,a4paper]{article}
\usepackage[margin=2.2cm]{geometry}
\usepackage{amsmath,amssymb}
\usepackage{booktabs}
\usepackage{enumitem}
\usepackage{xcolor}
\usepackage{hyperref}
\usepackage{natbib}
\usepackage{graphicx}
\usepackage{fancyhdr}
\usepackage{titlesec}
\usepackage{tcolorbox}

\definecolor{critred}{HTML}{D55E00}
\definecolor{warnorg}{HTML}{E69F00}
\definecolor{okblue}{HTML}{0072B2}
\definecolor{okgreen}{HTML}{009E73}

\newtcolorbox{critical}{colback=critred!8, colframe=critred!60, title={\textbf{Critical Issue}}, fonttitle=\bfseries\color{white}, coltitle=white, colbacktitle=critred!80}
\newtcolorbox{warning}{colback=warnorg!8, colframe=warnorg!60, title={\textbf{Warning}}, fonttitle=\bfseries\color{white}, coltitle=white, colbacktitle=warnorg!80}
\newtcolorbox{strength}{colback=okgreen!8, colframe=okgreen!60, title={\textbf{Strength}}, fonttitle=\bfseries\color{white}, coltitle=white, colbacktitle=okgreen!80}
\newtcolorbox{actionbox}[1][]{colback=okblue!8, colframe=okblue!60, title={#1}, fonttitle=\bfseries\color{white}, coltitle=white, colbacktitle=okblue!80}

\pagestyle{fancy}
\fancyhf{}
\rhead{\textit{Critical Review --- Europa Ice Shell Convection}}
\lhead{\textit{Pre-Thesis Assessment}}
\cfoot{\thepage}

\titleformat{\section}{\Large\bfseries\color{okblue}}{}{0em}{}
\titleformat{\subsection}{\large\bfseries}{}{0em}{}
\titleformat{\subsubsection}{\normalsize\bfseries\itshape}{}{0em}{}

\title{\textbf{Critical Scientific Review}\\[0.3em]
\Large Europa Ice Shell Convection Model:\\
Physics, Mathematics, Code, and Research Direction\\[0.5em]
\large\textit{Pre-Thesis Assessment \& Actionable Goals}}
\author{Generated for Joshua --- March 2026}
\date{}

\begin{document}
\maketitle
\thispagestyle{fancy}

\begin{abstract}
This document provides a systematic critical review of the EuropaConvection
modelling project ahead of thesis writing. It evaluates the physics model,
mathematical framework, Monte Carlo uncertainty quantification, Juno
observational constraining, and the central scientific hypothesis: that
ocean-driven heat delivery to Europa's equator can explain the observed
concentration of chaos terrain at low latitudes. The review identifies
\textbf{three critical gaps} that must be addressed and proposes
\textbf{three actionable goals} to close them before the thesis is written.
\end{abstract}

\tableofcontents
\newpage

%======================================================================
\section{The Central Hypothesis}
%======================================================================

\subsection{What you set out to show}

Europa's surface displays a striking geographic asymmetry: chaos terrain
(disrupted ice blocks, melt-through features like Conamara Chaos) clusters
at equatorial and tropical latitudes, while polar regions are dominated
by organised ridged terrain. This is counter-intuitive because tidal
heating is expected to peak at the poles ($\varepsilon_\mathrm{pole}
\approx 2\times \varepsilon_\mathrm{eq}$; Tobie et al.\ 2003, Beuthe 2013).

Your hypothesis:
\begin{quote}
\textit{Ocean circulation preferentially delivers heat to the equatorial
ice--ocean interface (Soderlund et al.\ 2014), thinning the equatorial ice
shell and promoting melt-through, diapirism, and surface disruption ---
producing the observed chaos terrain distribution.}
\end{quote}

\subsection{What the model actually tests}

The model tests whether \textbf{varying the local basal heat flux} at
equatorial surface conditions produces shell thicknesses that are
(a)~thinner at the equator than at the poles, and (b)~consistent with the
Juno MWR conductive-lid constraint of $D_\mathrm{cond} = 29 \pm 10$~km.
It does this via:

\begin{enumerate}[leftmargin=2em]
\item A 1D transient thermal solver with stagnant-lid convection
      parameterisation (Green et al.\ 2021).
\item Monte Carlo uncertainty propagation over 12+ physical parameters.
\item Bayesian importance reweighting against the Juno observation.
\item A 2D latitude-dependent extension with 37 axisymmetric columns.
\end{enumerate}

%======================================================================
\section{Review of the Physics Model}
%======================================================================

\subsection{Governing equations}

The model solves the 1D transient heat equation with volumetric tidal
heating:
\begin{equation}
\rho\, c_p \frac{\partial T}{\partial t}
  = \nabla \cdot (k \nabla T) + q_\mathrm{tidal}(T)
\label{eq:heat}
\end{equation}
discretised with Crank--Nicolson (after Rannacher startup) on 31 radial
nodes. Harmonic-mean conductivity averaging at interfaces ensures flux
conservation across the conductive--convective transition.

\begin{strength}
The numerical scheme is well-chosen: Rannacher startup avoids spurious
oscillations from stiff initial conditions, Picard iteration handles the
nonlinear $k(T)$ and $q_\mathrm{tidal}(T)$ couplings, and the tridiagonal
solve is $O(N)$. Auto-switching between spherical and Cartesian geometry
at $H = 30$~km is a sensible practical choice.
\end{strength}

\subsection{Material properties}

Temperature-dependent thermal conductivity $k(T) = 567/T$~W\,m$^{-1}$\,K$^{-1}$
(Howell 2021) is the \textbf{single most important equation in the model}.
It creates a factor-of-4 conductivity contrast between polar ($T \sim 50$~K,
$k \sim 11$~W/m/K) and equatorial ($T \sim 110$~K, $k \sim 5$~W/m/K) surfaces.

\begin{critical}
Your own diagnostic (2026-03-20 surface-temperature-dominance analysis)
showed that \textbf{surface temperature alone accounts for $\sim$101\% of the
pole--equator lid-thickness contrast}. This means the $k(T)$ relationship,
not ocean heat transport, is doing essentially all the work. The ocean
transport multiplier adds a measurable but \textbf{secondary} correction.
This finding needs to be front and centre in the thesis, not buried.
\end{critical}

\subsection{Rheology and viscosity}

The composite diffusion-creep viscosity follows Howell (2021):
\begin{equation}
\eta = \frac{1}{2}\left[
  \frac{42\, V_m}{R\, T\, d^2}
  \left(D_v + \frac{\pi\delta}{d}\,D_b\right)
\right]^{-1}
\label{eq:visc}
\end{equation}
with volume and grain-boundary diffusivities $D_v, D_b$ following Arrhenius
laws.

\begin{strength}
Supporting both Maxwell and Andrade rheologies for tidal dissipation is
good practice. The Andrade formulation follows McCarthy et al.\ (2011) and
Renaud \& Henning (2018), which is the current community standard for icy
satellites.
\end{strength}

\begin{warning}
The grain-size transition at $T = 230$~K (1.0~mm $\to$ 0.1~mm) is a
hard step function. In reality, dynamic recrystallisation produces a
continuous grain-size profile that depends on strain rate, temperature,
and stress history (Barr \& McKinnon 2007). This step may introduce
an artificial viscosity jump at 230~K.
\end{warning}

\subsection{Stagnant-lid convection parameterisation}

The model uses the Green et al.\ (2021) / Deschamps \& Vilella (2021)
temperature-profile scanning method:
\begin{equation}
\mathrm{Nu} = C\,\mathrm{Ra}^\xi
  \left(\frac{T_i - T_c}{\Delta T}\right)^\zeta, \qquad
\mathrm{Ra}_\mathrm{crit} = 1000
\end{equation}

\begin{strength}
Dynamic interface detection (scanning for $T \geq T_c$ in the profile)
is physically more defensible than fixed-depth cutoffs. The convection
ramp for cold columns prevents numerical discontinuities.
\end{strength}

%======================================================================
\section{Review of the Monte Carlo Framework}
%======================================================================

\subsection{Prior distributions}

The 2026 parameter audit is well-documented and literature-grounded:

\begin{center}
\begin{tabular}{llll}
\toprule
\textbf{Parameter} & \textbf{Prior} & \textbf{Range} & \textbf{Sensitivity} \\
\midrule
$d_\mathrm{grain}$ & LogNormal(0.6~mm, $\sigma=0.35$~dex) & [0.05, 3.0]~mm & \textbf{High} \\
$\varepsilon_0$ & LogNormal($1.2\times10^{-5}$, $\sigma=0.3$) & [$2\times10^{-6}$, $3.4\times10^{-5}$] & \textbf{High} \\
$\mu_\mathrm{ice}$ & Normal($3.5\pm0.5$~GPa) & [2, 5]~GPa & High \\
$q_\mathrm{basal}$ & Uniform & [5, 25]~mW/m$^2$ & Medium \\
$T_\mathrm{surf}$ & Normal($104\pm7$~K) & [80, 120]~K & Medium \\
$Q_v$ & Normal($59.4$~kJ/mol, 5\%) & --- & \textbf{Dominant} \\
\bottomrule
\end{tabular}
\end{center}

\begin{strength}
Replacing the old $P_\mathrm{tidal}$ sampling with direct $q_\mathrm{basal}$
is a significant improvement. Fixing $f_\mathrm{salt} = 0$ (pure-ice
baseline) with salinity as a separate scenario avoids the unconstrained
salt--conductivity degeneracy of the old prior.
\end{strength}

\subsection{Sobol sensitivity analysis}

The Sobol decomposition reveals the model's structural truth:

\begin{center}
\begin{tabular}{lcc}
\toprule
\textbf{Parameter} & $S_1$ (first-order) & $S_T$ (total-order) \\
\midrule
$Q_v$ (activation energy) & 0.365 & \textbf{0.557} \\
$d_\mathrm{grain}$ (grain size) & 0.313 & \textbf{0.496} \\
$\varepsilon_0$ (tidal strain) & 0.096 & 0.285 \\
$q_\mathrm{basal}$ (ocean heat) & 0.044 & \textbf{0.142} \\
$T_\mathrm{surf}$ (surface temp) & 0.014 & 0.023 \\
\bottomrule
\end{tabular}
\end{center}

\begin{critical}
\textbf{Rheology dominates the conductive lid.} The combined $S_T$ for
$Q_v + d_\mathrm{grain} = 1.05$, while $q_\mathrm{basal}$ contributes
only $S_T = 0.14$. This means the ocean heat flux --- the physical quantity
your hypothesis depends on --- explains only $\sim$14\% of the variance in
$D_\mathrm{cond}$. The model's own sensitivity analysis shows that
\textbf{the observable you compare to Juno ($D_\mathrm{cond}$) is largely
insensitive to the mechanism you are testing (ocean heat transport)}.

This is the most important finding in the project and must be honestly
confronted in the thesis.
\end{critical}

%======================================================================
\section{Review of Juno Constraining}
%======================================================================

\subsection{Methodology}

Bayesian importance reweighting against $D_\mathrm{cond}^\mathrm{obs}
= 29 \pm 10$~km (Levin et al.\ 2026) with iterative prior tightening.
Converged in a single round for all three ocean scenarios at 35$^\circ$
latitude, with ESS fractions $\sim$69\%.

\begin{strength}
Importance reweighting is the right choice for an expensive forward model.
The ESS fraction of 69\% indicates good prior--likelihood overlap, and
the PSIS $\hat{k} \approx 0.45$ confirms stable tail behaviour.
\end{strength}

\subsection{Results}

All three ocean transport models produce nearly identical posteriors at
35$^\circ$:

\begin{center}
\begin{tabular}{lccc}
\toprule
\textbf{Scenario} & $q_\mathrm{mult}$ & $D_\mathrm{cond}$ posterior & ESS/N \\
\midrule
Uniform & 1.000 & 22.24~km & 69.1\% \\
Soderlund & 1.002 & 22.23~km & 69.1\% \\
Lemasquerier & 0.998 & 22.26~km & 69.1\% \\
\bottomrule
\end{tabular}
\end{center}

\begin{critical}
\textbf{The Juno data cannot discriminate between ocean transport modes.}
At 35$^\circ$ latitude, the transport multipliers are all $\approx 1.0$,
so the three scenarios collapse. This is not a failure of the method but a
geometric limitation: a single mid-latitude measurement point has no
leverage on equator--pole contrasts. The equatorial Bayesian refit
(posterior\_summary.csv) confirms this: log Bayes factors between scenarios
are all $<0.3$, far below the threshold for ``substantial'' evidence.

The thesis must clearly state: \textit{Juno's mid-latitude observation
constrains the global mean shell properties but cannot resolve the
latitudinal heat transport pattern.}
\end{critical}

\subsection{What the Juno refit did constrain}

The meaningful posterior updates are in rheological parameters:
\begin{itemize}[leftmargin=2em]
\item $d_\mathrm{grain}$: shifted from 0.60~mm (prior) to 0.76~mm
      (posterior), 90\% CI narrowed by 55\% (from 1.78 to 0.81~dex).
      But 0.31--1.96~mm is still nearly an order of magnitude.
\item $q_\mathrm{basal}$: 90\% CI narrowed by 19\% (from 20 to 16~mW/m$^2$
      width). Median shifted to 10.3~mW/m$^2$ from prior centre of 15.
\end{itemize}

\begin{warning}
The posterior $d_\mathrm{grain}$ range of [0.31, 1.96]~mm is still
broad enough that it cannot distinguish between the Mitri (2023) estimate
of 0.39--0.80~mm and the Ruiz et al.\ (2007) range of 0.2--2~mm.
\end{warning}

%======================================================================
\section{Review of the 2D Latitude Model}
%======================================================================

\subsection{Architecture}

The 2D model couples 37 independent 1D columns via operator splitting:
radial heat equation first, then explicit lateral diffusion between columns.

\begin{warning}
\textbf{Lateral diffusion is negligible.} The thermal diffusion timescale
across Europa's radius is $\tau \sim R^2/\kappa \sim 200$~Gyr, far exceeding
the model's 15.8~Myr integration window. The 2D model therefore reduces to
37 independent 1D columns with latitude-varying boundary conditions. This is
not wrong, but the ``2D'' framing may overstate the physics. The lateral
coupling that matters is \textbf{ocean circulation}, which is imposed as a
boundary condition, not solved.
\end{warning}

\subsection{Latitude-dependent physics}

The latitude profile includes:
\begin{enumerate}[leftmargin=2em]
\item Surface temperature: $T_s(\phi) = [(T_\mathrm{eq}^4 - T_\mathrm{floor}^4)
      \cos^p\phi + T_\mathrm{floor}^4]^{1/4}$, calibrated to Ashkenazy (2019).
\item Tidal strain: $\varepsilon_0(\phi) = \varepsilon_\mathrm{eq}
      \sqrt{1 + c\sin^2\phi}$ (Beuthe 2013 whole-shell pattern).
\item Ocean heat flux: scenario-dependent pattern $q(\phi)$ with three options.
\end{enumerate}

\begin{strength}
The surface temperature calibration ($p = 1.25$, RMS~$= 0.70$~K vs
seasonal energy balance) is well-validated. The Beuthe (2013) strain
patterns with multiple geometries (mantle-core, shell-dominated,
non-monotonic) provide good coverage of the tidal dissipation uncertainty.
\end{strength}

%======================================================================
\section{Honest Assessment: The Hypothesis Gap}
%======================================================================

Here is the uncomfortable truth that the thesis must confront head-on:

\subsection{The chain of inference has a missing link}

Your hypothesis chain is:
\[
\underbrace{\text{Ocean heat}}_{\text{cause}}
\;\longrightarrow\;
\underbrace{\text{Thin equatorial ice}}_{\text{mechanism}}
\;\longrightarrow\;
\underbrace{\text{Chaos terrain}}_{\text{observation}}
\]

Your model tests the first arrow ($\longrightarrow$) but \textbf{never
tests the second}. You show that enhanced equatorial heat flux thins
$H_\mathrm{total}$ from 35~km to 24~km (for 1.5$\times$ enhancement), but
you do not model:
\begin{itemize}[leftmargin=2em]
\item What shell thickness triggers melt-through or diapir formation.
\item Whether thinner ice at the equator is mechanically weaker.
\item Whether the predicted thickness contrast matches the observed chaos
      distribution.
\end{itemize}

\subsection{$D_\mathrm{cond}$ is the wrong observable for this hypothesis}

The Juno MWR measures the conductive lid thickness $D_\mathrm{cond}$. But
your Sobol analysis shows $D_\mathrm{cond}$ is rheology-dominated
($S_T = 1.05$ for rheological parameters). Meanwhile, $H_\mathrm{total}$
--- the mechanically relevant quantity for melt-through --- \textit{is}
sensitive to $q_\mathrm{basal}$ (varying 23--51~km across transport
scenarios). The model predicts a strong signal in $H_\mathrm{total}$ but
Juno only measures $D_\mathrm{cond}$, where the signal is weak.

\begin{critical}
\textbf{The observable you have (Juno $D_\mathrm{cond}$) is insensitive to
the mechanism you care about (ocean heat). The quantity that is sensitive
($H_\mathrm{total}$) is not directly observable.} This is not a flaw in
your work --- it is a fundamental limitation of the current observational
constraints that the thesis should frame honestly.
\end{critical}

\subsection{Surface temperature dominance}

Your own analysis (2026-03-20) showed that surface temperature alone
produces the equator--pole lid contrast. Ocean heat transport adds a
correction, but the baseline story is:
\begin{quote}
\textit{Equatorial ice has a thinner conductive lid because warm equatorial
surfaces reduce conductivity ($k = 567/T$), not primarily because the ocean
delivers more heat there.}
\end{quote}

This is actually a \textbf{robust, defensible result}. It just shifts the
narrative from ``ocean dynamics cause chaos'' to ``surface thermal physics
creates the latitude structure, and ocean heat modulates it.''

%======================================================================
\section{What You Should Actually Be Doing}
%======================================================================

Given the model's strengths and the gaps identified above, here are the
three areas where your effort will have the highest return before thesis
writing.

\subsection{The question you can answer}

You cannot definitively prove that ocean heat causes chaos terrain with this
model. But you \textit{can} answer a well-posed, publishable question:

\begin{actionbox}[\textbf{Thesis Question}]
\textit{What is the predicted latitude-dependent ice shell structure on
Europa under competing ocean heat transport scenarios, and which scenarios
are consistent with the Juno MWR constraint?}
\end{actionbox}

This reframes the work from ``explaining chaos'' to ``predicting shell
structure'' --- a question your model is well-designed to answer.

%======================================================================
\section{Three Actionable Goals Before Thesis Writing}
%======================================================================

\begin{actionbox}[\textbf{Goal 1: Produce the definitive pole-to-equator shell profiles}]
\textbf{What}: Run the 2D latitude model (37 columns, 0$^\circ$--90$^\circ$)
for $N \geq 500$ MC realisations under all three ocean scenarios (uniform,
Soderlund, Lemasquerier), with Andrade rheology and the audited 2026 priors.
Produce latitude profiles of:
\begin{itemize}[leftmargin=1.5em]
\item $D_\mathrm{cond}(\phi)$, $D_\mathrm{conv}(\phi)$, $H_\mathrm{total}(\phi)$
      (median $\pm$ 1$\sigma$ envelopes)
\item Convective fraction and Ra/Nu as functions of latitude
\item Pole-to-equator thickness ratio $\Delta H / H_\mathrm{mean}$
\end{itemize}

\textbf{Why}: This is the \textbf{core deliverable} of the thesis. It
provides the predicted observables that Europa Clipper can test.
Currently you have 2D runs at $N = 5$--500 but no systematic comparison
figure across all three scenarios at production sample size with
uncertainty envelopes.

\textbf{Deliverable}: One publication-quality multi-panel figure showing
$H_\mathrm{total}(\phi)$ and $D_\mathrm{cond}(\phi)$ for three scenarios
with uncertainty bands. This is Figure~1 of the thesis.

\textbf{Time estimate}: 2--3 days (MC runs + figure generation).
\end{actionbox}

\vspace{0.5em}

\begin{actionbox}[\textbf{Goal 2: Honest Juno comparison with model discrimination}]
\textbf{What}: Apply the Bayesian importance reweighting to the 2D
latitude profiles (not just 1D endmembers), using the Juno $D_\mathrm{cond}$
constraint at 35$^\circ$. Compute:
\begin{itemize}[leftmargin=1.5em]
\item Posterior $D_\mathrm{cond}$ at 35$^\circ$ for each scenario
\item Bayes factors between scenarios (uniform vs Soderlund vs Lemasquerier)
\item Posterior predictive distributions for $D_\mathrm{cond}$ at the equator
      and pole (predictions for Clipper)
\end{itemize}

\textbf{Why}: The current Juno constraining is done on 1D mid-latitude
columns, which collapse all scenarios to identical posteriors. By using the
full 2D profiles, you can ask: ``Given that Juno sees $\sim$29~km at 35$^\circ$,
what does each scenario predict at 0$^\circ$ and 90$^\circ$?'' This is where
scenario discrimination appears --- not at the observation latitude, but
at the prediction latitudes.

\textbf{Deliverable}: A table of scenario-conditional predictions at
equator/pole/mid-latitude. A figure showing posterior predictive
$D_\mathrm{cond}(\phi)$ under each scenario. Bayes factors reported
honestly (likely inconclusive, which is itself a result).

\textbf{Time estimate}: 3--4 days (reweighting + figure + interpretation).
\end{actionbox}

\vspace{0.5em}

\begin{actionbox}[\textbf{Goal 3: Sensitivity decomposition by physical mechanism}]
\textbf{What}: Decompose the pole-to-equator thickness contrast into
contributions from:
\begin{enumerate}[leftmargin=1.5em]
\item Surface temperature ($k(T)$ effect alone)
\item Tidal heating latitude pattern ($\varepsilon_0(\phi)$ effect)
\item Ocean heat transport ($q_\mathrm{basal}(\phi)$ effect)
\item Interaction terms
\end{enumerate}

Do this by running the 2D model with each effect turned on/off:
\begin{itemize}[leftmargin=1.5em]
\item Run A: uniform $T_s$, uniform $\varepsilon_0$, uniform $q$ (control)
\item Run B: latitude-varying $T_s$ only (isolate surface effect)
\item Run C: latitude-varying $T_s + \varepsilon_0$ (add tidal)
\item Run D: latitude-varying all (full model, per scenario)
\end{itemize}

\textbf{Why}: Your 2026-03-20 surface-temperature-dominance analysis is
the most important finding in the project, but it was done as a one-off
diagnostic, not as a systematic decomposition with uncertainty. Elevating
this to a proper factor-separation experiment (Stein \& Alpert 1993) gives
you a defensible, quantitative attribution of what controls Europa's shell
structure. This is the intellectual core of the thesis.

\textbf{Deliverable}: A bar chart or waterfall diagram showing the
contribution of each physical mechanism to $\Delta H_\mathrm{total}$
(pole $-$ equator) with uncertainty bands. This is the ``money figure'' ---
the one that says ``surface temperature does X\%, tidal does Y\%, ocean
does Z\%.''

\textbf{Time estimate}: 3--4 days (4 run sets $\times$ 3 scenarios + analysis).
\end{actionbox}

%======================================================================
\section{Thesis Narrative Recommendation}
%======================================================================

Based on this review, the honest thesis narrative is:

\begin{enumerate}[leftmargin=2em]
\item \textbf{Europa's ice shell structure is primarily controlled by
      rheology and surface temperature}, not by ocean heat transport.
      The temperature-dependent conductivity $k(T) = 567/T$ creates a
      factor-of-4 conductivity contrast between poles and equator, which
      alone produces the observed lid-thickness latitude gradient.

\item \textbf{Ocean heat transport modulates but does not dominate the
      shell structure.} Under Soderlund-type equatorial enhancement, the
      total shell thickness changes by $\sim$12~km (equatorial thinning),
      but the conductive lid --- the Juno-observable quantity --- changes
      by $<$2~km.

\item \textbf{The Juno MWR constraint is consistent with all tested
      ocean scenarios} because the observed quantity ($D_\mathrm{cond}$)
      is rheology-dominated. Discriminating between scenarios requires
      either (a)~equator--pole $D_\mathrm{cond}$ measurements (Clipper),
      or (b)~total shell thickness observations (radar, gravity).

\item \textbf{The chaos terrain connection remains an open question.}
      The model predicts equatorial thinning consistent with enhanced
      melt probability, but does not model the melt-through process itself.
      This is a prediction for future coupled thermomechanical models.
\end{enumerate}

This narrative is strong, honest, and publishable. It doesn't overstate the
ocean-heat story but frames it as a \textit{testable prediction} rather than
a proven mechanism.

%======================================================================
\section{Summary of Critical Issues}
%======================================================================

\begin{center}
\begin{tabular}{p{0.5cm}p{5.5cm}p{3cm}p{4cm}}
\toprule
\textbf{\#} & \textbf{Issue} & \textbf{Severity} & \textbf{Resolution} \\
\midrule
1 & $D_\mathrm{cond}$ is rheology-dominated, not ocean-heat-driven
  & \textcolor{critred}{\textbf{Critical}} & Reframe thesis; use $H_\mathrm{total}$ \\
2 & Juno cannot discriminate ocean scenarios at 35$^\circ$
  & \textcolor{critred}{\textbf{Critical}} & Predict at other latitudes \\
3 & Chaos terrain link untested (no melt-through model)
  & \textcolor{warnorg}{\textbf{Major}} & Acknowledge as prediction \\
4 & 2D lateral diffusion is negligible ($\tau \sim 200$~Gyr)
  & \textcolor{warnorg}{\textbf{Moderate}} & Reframe as 1D$\times$37 \\
5 & Grain-size step function at 230~K
  & \textcolor{okblue}{Minor} & Sensitivity test \\
6 & Only Andrade rheology in production
  & \textcolor{okblue}{Minor} & Add Maxwell comparison \\
\bottomrule
\end{tabular}
\end{center}

%======================================================================
\section{Summary of Strengths}
%======================================================================

\begin{enumerate}[leftmargin=2em]
\item \textbf{Rigorous parameter audit} with literature-grounded priors
      for all 12+ parameters.
\item \textbf{Sobol sensitivity analysis} providing quantitative attribution
      of variance to individual parameters --- rare in planetary science
      thermal modelling.
\item \textbf{Bayesian observational constraining} via importance reweighting
      with proper diagnostic checks (ESS, PSIS $\hat{k}$).
\item \textbf{Three competing ocean scenarios} spanning the published
      range of ocean heat transport patterns.
\item \textbf{Surface temperature calibration} validated against seasonal
      energy balance models.
\item \textbf{Reproducible codebase} with audited samplers, version-controlled
      results, and automated figure generation.
\end{enumerate}

%======================================================================
\section{Appendix: Key Equations Reference}
%======================================================================

\subsection{Thermal conductivity}
\begin{equation}
k(T) = \frac{567}{T} \quad [\mathrm{W\,m^{-1}\,K^{-1}}]
\end{equation}

\subsection{Composite viscosity (Howell 2021)}
\begin{equation}
\eta = \frac{1}{2}\left[
  \frac{42\, V_m}{R\, T\, d^2}
  \left(
    D_{0v}\,e^{-Q_v/RT}
    + \frac{\pi\delta}{d}\,D_{0b}\,e^{-Q_b/RT}
  \right)
\right]^{-1}
\end{equation}

\subsection{Maxwell tidal heating}
\begin{equation}
q_\mathrm{tidal} = \frac{\varepsilon_0^2\,\omega^2\,\eta}
  {2\left(1 + \omega^2\eta^2/\mu^2\right)}
\end{equation}

\subsection{Andrade complex compliance}
\begin{equation}
J^*(\omega) = J_U + \beta\,\Gamma(1+\alpha)\,(\omega\tau)^{-\alpha}\,e^{-i\alpha\pi/2}
  - \frac{i}{\omega\eta}
\end{equation}

\subsection{Rayleigh number}
\begin{equation}
\mathrm{Ra} = \frac{\rho\,g\,\alpha_T\,\Delta T\,D_\mathrm{conv}^3}{\kappa\,\eta}
\end{equation}

\subsection{Surface temperature profile}
\begin{equation}
T_s(\phi) = \left[
  (T_\mathrm{eq}^4 - T_\mathrm{floor}^4)\cos^p\phi + T_\mathrm{floor}^4
\right]^{1/4}, \quad p = 1.25
\end{equation}

\subsection{Tidal strain latitude dependence (Beuthe 2013)}
\begin{equation}
\varepsilon_0(\phi) = \varepsilon_\mathrm{eq}\,\sqrt{1 + c\,\sin^2\phi},
\quad c = \left(\frac{\varepsilon_\mathrm{pole}}{\varepsilon_\mathrm{eq}}\right)^2 - 1
\end{equation}

\end{document}
